\section{Introduction}

In the era of information overload, social media platforms have become pivotal sources of real-time content, offering vast amounts of data that can be both overwhelming and challenging to interpret. To address this issue, we developed an Automated Journalist App, designed to streamline and distill social media feeds into meaningful summaries.

Our project aims to enhance the accessibility and comprehensibility of social media content by implementing advanced summarization techniques. The application retrieves data from various social media platforms and employs sophisticated natural language processing methods to generate concise and relevant summaries. This process includes two key summarization approaches: topic-aware summarization and delta summarization.

Topic-aware summarization involves generating summaries of the same content from multiple perspectives or topics, allowing users to gain insights into different aspects of the information presented. Delta summarization, on the other hand, facilitates a comparative analysis of summaries across different topics, highlighting the variations and similarities between them.

By integrating these techniques, our application not only provides users with streamlined content but also offers a nuanced understanding of the information through comparative analysis. This work aims to contribute to the field of automated content summarization by enhancing the utility and interpretability of social media data.

The Automated Journalist App is a group work of the authors of this paper. In the project, Can Bilgin was mainly responsible for selecting the technology stack, code quality, and software architecture. Furthermore, Can Bilgin implemented the communication between third-party social media platforms and our Automated Journalist App. Ecem Varma and Ahmet Bilal Akın worked together to conduct research and made decisions about which NLP techniques should be employed in this project. Alongside the research task, Ecem Varma had an active role in implementing the front end and Ahmet Bilal Akın was responsible for implementing the NLP techniques into the backend. However, this does not mean that we did distinct tasks and each person was only responsible for his/her section. Throughout the whole project, we exchanged our knowledge about the related topics and helped each other out while implementing our parts.  